\begin{proposition}[投影自能的解析性、能量变化与开放通道]
设 $H=H^\dagger$，$P=P^2=P^\dagger$，$Q=I-P$，并在所研究工作空间内令
$B=QHP$ 有界，$H_Q=QHQ$ 自伴。下列逆算符均取在对应子空间上。
若 $z\in\rho(H_Q)$ 且 $z-PHP-B^\dagger(z-H_Q)^{-1}B$ 可逆，则
\[
P(z-H)^{-1}P=[z-PHP-\Sigma_P(z)]^{-1},
\qquad \Sigma_P(z)=B^\dagger(z-H_Q)^{-1}B.
\]
若实能量区间 $\mathcal W_E$ 与 $\sigma(H_Q)$ 相距 $d_Q>0$，则
\begin{align*}
\|\Sigma_P(E)\|&\le\frac{\|B\|^2}{d_Q},\\
\partial_E\Sigma_P(E)&=-B^\dagger(E-H_Q)^{-2}B,\\
\|\Sigma_P(E)-\Sigma_P(E_*)\|&\le
\frac{\|B\|^2}{d_Q^2}|E-E_*|.
\end{align*}
最后一个界要求连接 $E,E_*$ 的整段区间均满足同一谱隙条件。
\end{proposition}
\begin{proof}
对源 $f=Pf$ 写 $\Psi=p+q$，其中 $p=P\Psi,q=Q\Psi$。
分块方程是
\[
(z-PHP)p-B^\dagger q=f,\qquad (z-H_Q)q=Bp.
\]
由第二式得 $q=(z-H_Q)^{-1}Bp$；代回第一式即得到投影预解式。
注意逆算符求导带负号：
\[
\frac{\dd}{\dd E}(E-H_Q)^{-1}=-(E-H_Q)^{-2}.
\]
自伴谱定理给 $\|(E-H_Q)^{-1}\|\le d_Q^{-1}$。
对导数界在 $E_*$ 至 $E$ 积分，便得到自能变化界。
因此有隙闭合通道中的自能沿实能量单调不增；它的频率依赖不能任意指定。
\end{proof}

\begin{proposition}[谱密度与自能的损失号]
令 $\mathsf P_Q(\dd\varepsilon)$ 为 $H_Q$ 的谱投影测度。
则 $\mathsf J_P(\dd\varepsilon)=B^\dagger\mathsf P_Q(\dd\varepsilon)B$
是取值于正算符的测度，且
\begin{equation}
\Sigma_P(z)=\int\frac{\mathsf J_P(\dd\varepsilon)}{z-\varepsilon}.
\label{eq:feqo-feshbach-positive-spectral-measure}
\end{equation}
在绝对连续谱区域，若密度 $J_P(E)$ 的边界值存在，则
\begin{equation}
\Sigma_P(E+\ii0)=\mathop{\rm PV}\int
\frac{J_P(\varepsilon)}{E-\varepsilon}\,\dd\varepsilon
-\ii\pi J_P(E),\qquad
\Gamma_P(E)=\frac{2\pi}{\hbar}J_P(E)\succeq0.
\label{eq:feqo-feshbach-width-positive}
\end{equation}
这里 $J_P$ 的单位为能量，$\Gamma_P$ 的单位为 $\si{\per\second}$。
\end{proposition}
\begin{proof}
将 $H_Q$ 的谱积分代入预解式即可得到第一式。
对任意保留态 $u$，$\langle u|\mathsf J_P(\Delta)|u\rangle
=\|\mathsf P_Q(\Delta)Bu\|^2\ge0$，证明正性。
再用 $(x+\ii0)^{-1}=\mathop{\rm PV}(1/x)-\ii\pi\delta(x)$ 得到边界值。
因而衰减极点的负虚部来自正谱权重。对隔离的标量共振，把其近似为常数时，振幅含
$\exp(-\Gamma_Pt/2)$，概率含 $\exp(-\Gamma_Pt)$。
一般矩阵情形的有效 Hamiltonian 是 $H_{\rm ren}-\ii\hbar\Gamma_P/2$；
若厄米部分与 $\Gamma_P$ 不对易，不能把整个传播子拆成独立的幺正因子和这个衰减因子。
此时被消去通道仍携带相应损失概率；仅传播这个非厄米保留振幅不能得到保迹量子信道。
\end{proof}

同一消元还能直接在时域看见记忆。若 $q(0)$ 给定，解第二个 Schr\"odinger 分块方程得到
\[
q(t)=\ee^{-\ii H_Qt/\hbar}q(0)
-\frac{\ii}{\hbar}\int_0^t\ee^{-\ii H_Q(t-s)/\hbar}Bp(s)\,\dd s.
\]
因此保留振幅满足
\begin{equation}
\ii\hbar\dot p(t)=PHPp(t)+B^\dagger\ee^{-\ii H_Qt/\hbar}q(0)
-\frac{\ii}{\hbar}\int_0^t B^\dagger\ee^{-\ii H_Q(t-s)/\hbar}Bp(s)\,\dd s.
\label{eq:feqo-feshbach-exact-amplitude-memory}
\end{equation}
频域自能与时域记忆核由同一个谱测度决定。
初始 $Q$ 分量产生独立源项，谱隙界控制闭合通道，连续谱边界值描述真实开放通道；
这三件事将在开放系统的投影方程中以相同的分块代数再次出现。
